\documentclass[border=5pt]{standalone}
\usepackage{pgfplots}
\usepackage{pgf-pie} % 饼图：AI 代码可直接使用 \pie
\pgfplotsset{compat=1.18}
\usepgfplotslibrary{fillbetween}  % 面积图 / 堆叠面积图 / \closedcycle 填充路径
% 注：error bars 是 pgfplots 内置功能（在核心 pgfplots.errorbars.code.tex），不需要单独 \usepgfplotslibrary 加载
\usepackage{amsmath}
\usepackage{amssymb}
\usepackage{xcolor}[dvipsnames,svgnames]  % 加载标准色名表（steelblue/teal/orange/coral 等），避免 AI 用预定义色名时报 Undefined color

% 支持中文
\usepackage{fontspec}
\usepackage{xeCJK}
\setCJKmainfont{SimSun}
\setmainfont{Times New Roman}

\begin{document}

\begin{tikzpicture}\begin{axis}[\n    ybar, width=14cm, height=7cm,\n    title={2023年北京各月降水量},\n    ylabel={降水量（毫米）},\n    symbolic x coords={1月,2月,3月,4月,5月,6月,7月,8月,9月,10月,11月,12月},\n    xtick=data,\n    x tick label style={font=\scriptsize},\n    bar width=9pt,\n    ymin=0, ymax=260,\n    enlarge x limits=0.15,\n    nodes near coords,\n    point meta=explicit symbolic,\n    every node near coord/.append style={anchor=south, font=\scriptsize, fill=none, draw=none, inner sep=1pt}\n]\n\addplot[fill=blue!60, draw=blue!60] coordinates {\n    (1月,1.4)[1.4]\n    (2月,5.6)[5.6]\n    (3月,8.3)[8.3]\n    (4月,21.5)[22]\n    (5月,34.2)[34]\n    (6月,78.4)[78]\n    (7月,221.0)[221]\n    (8月,165.3)[165]\n    (9月,42.6)[43]\n    (10月,18.5)[19]\n    (11月,12.3)[12]\n    (12月,2.1)[2.1]\n};\n\node[anchor=north west, font=\scriptsize] at (axis description cs:0.0,-0.20) {数据来源：北京市气象局 2023 年逐月降水量统计（约数）};\n\end{axis}\n\end{tikzpicture}

\end{document}